%============================================================
%  Bd. 24 - Idempotente Halbgruppen
%============================================================

\documentclass[main.tex]{subfiles}

\ifSubfilesClassLoaded{
  \BandSetupStandalone{B24}
}{
}

\title{Bd. 24 - Idempotente Halbgruppen}
\author{Martin Kunze}
\date{}
\setcounter{file}{24}

\hypersetup{
  pdftitle={Bd. 24 - Idempotente Halbgruppen},
  pdfauthor={Martin Kunze}
}

\begin{document}

\title{Bd. 24 - Idempotente Halbgruppen}
\author{Martin Kunze}
\date{}
\hypersetup{pageanchor=false}
\maketitle
\hypersetup{pageanchor=true}
\setcounter{tocdepth}{3}
\tableofcontents
\setcounter{file}{24}
\renewcommand{\FormulaBandID}{24}

\chapter{Einführung}

Dieser Band untersucht Halbgruppen mit idempotenter Operation.

\chapter{Idempotente Halbgruppen}

\section{Definition}

\FormulaDefDeltaK[Begriff der idempotenten Halbgruppe]{%
  \IdempotentSemigroupAlg{A}{\star}%
}{IdempotentSemigroupDef}{%
  \DeltaRow{Mengen}{A\dsep x}
  \DeltaRow{\textbf{Axiome}}{}
  \DeltaPrem{Halbgruppe}{\SemigroupAlg{A}{\star}}
  \DeltaRow{Idempotenz}{x\in A\vdash x\star x=x}
  \DeltaRow{\textbf{Neues Symbol}}{}
  \DeltaPrem{Idempotente Halbgruppe}{%
    \IdempotentSemigroupAlg{A}{\star}}
}

\FormulaAxiomDeltaK[Zugriff auf die Halbgruppe einer idempotenten Halbgruppe]{%
  \IdempotentSemigroupAlg{A}{\star}
  \vdash\SemigroupAlg{A}{\star}%
}{IdempotentSemigroupBaseAxiom}{%
  \DeltaRow{Mengen}{A}
}

\FormulaAxiomDeltaK[Zugriff auf die Idempotenz]{%
  \IdempotentSemigroupAlg{A}{\star}\dsep x\in A
  \vdash x\star x=x%
}{IdempotentSemigroupIdempotenceAxiom}{%
  \DeltaRow{Mengen}{A\dsep x}
}

\section{Beispiele}

\FormulaThmDeltaK[Die Vereinigung bildet auf einer Potenzmenge eine
idempotente Halbgruppe]{%
  \IdempotentSemigroupAlg{\powerset(M)}{\cup}%
}{PowerSetUnionIdempotentSemigroup}{%
  \DeltaRow{Mengen}{M\dsep X}
}
\begin{tabproof}
  \proofstep{}{\CommutativeSemigroupAlg{\powerset(M)}{\cup}}{%
    \FormulaRefAuto{PowerSetUnionCommutativeSemigroup}[thm]}
  \proofstep{}{\SemigroupAlg{\powerset(M)}{\cup}}{%
    \FormulaRefAuto{CommutativeSemigroupBaseAxiom}[ax]{1}}
  \proofstep{3}{X\in\powerset(M)}{\rA}
  \proofstep{}{X\cup X=X}{%
    \FormulaRefAuto{A \cup A = A}}
  \proofstep{}{\IdempotentSemigroupAlg{\powerset(M)}{\cup}}{%
    \FormulaRefAuto{IdempotentSemigroupDef}[def]{2,4}}
\end{tabproof}

\paragraph{Nichtkommutatives Beispiel.}
Für eine Menge $A$ mit mindestens zwei Elementen sei
$x\star_{\ell}y\coloneqq x$. Dann gelten
$(x\star_{\ell}y)\star_{\ell}z=x
=x\star_{\ell}(y\star_{\ell}z)$ und
$x\star_{\ell}x=x$, aber im Allgemeinen
$x\star_{\ell}y\neq y\star_{\ell}x$.

\section{Morphismen}

\FormulaDefDeltaK[Homomorphismus idempotenter Halbgruppen]{%
  \IdempotentSemigroupHom{f}{A,\star}{B,\diamond}%
}{IdempotentSemigroupHomDef}{%
  \DeltaRow{Mengen}{A\dsep B}
  \DeltaRow{Funktionen}{f\dsep\star\dsep\diamond}
  \DeltaRow{\textbf{Axiome}}{}
  \DeltaPrem{Quellstruktur}{\IdempotentSemigroupAlg{A}{\star}}
  \DeltaPrem{Zielstruktur}{\IdempotentSemigroupAlg{B}{\diamond}}
  \DeltaPrem{Halbgruppenhomomorphismus}{%
    \SemigroupHom{f}{A,\star}{B,\diamond}}
  \DeltaRow{\textbf{Neues Symbol}}{}
  \DeltaPrem{Homomorphismus}{%
    \IdempotentSemigroupHom{f}{A,\star}{B,\diamond}}
}

\FormulaAxiomDeltaK[Quellstruktur]{%
  \IdempotentSemigroupHom{f}{A,\star}{B,\diamond}
  \vdash\IdempotentSemigroupAlg{A}{\star}%
}{IdempotentSemigroupHomSourceAxiom}{%
  \DeltaRow{Mengen}{A\dsep B}
  \DeltaRow{Funktionen}{f\dsep\star\dsep\diamond}
}

\FormulaAxiomDeltaK[Zielstruktur]{%
  \IdempotentSemigroupHom{f}{A,\star}{B,\diamond}
  \vdash\IdempotentSemigroupAlg{B}{\diamond}%
}{IdempotentSemigroupHomTargetAxiom}{%
  \DeltaRow{Mengen}{A\dsep B}
  \DeltaRow{Funktionen}{f\dsep\star\dsep\diamond}
}

\FormulaAxiomDeltaK[Zugriff auf den Halbgruppenhomomorphismus]{%
  \IdempotentSemigroupHom{f}{A,\star}{B,\diamond}
  \vdash\SemigroupHom{f}{A,\star}{B,\diamond}%
}{IdempotentSemigroupHomBaseHomAxiom}{%
  \DeltaRow{Mengen}{A\dsep B}
  \DeltaRow{Funktionen}{f\dsep\star\dsep\diamond}
}

\FormulaDefDeltaK[Isomorphismus idempotenter Halbgruppen]{%
  \IdempotentSemigroupIso{f}{A,\star}{B,\diamond}%
}{IdempotentSemigroupIsoDef}{%
  \DeltaRow{Mengen}{A\dsep B}
  \DeltaRow{Funktionen}{f\dsep\star\dsep\diamond}
  \DeltaRow{\textbf{Axiome}}{}
  \DeltaPrem{Homomorphismus}{%
    \IdempotentSemigroupHom{f}{A,\star}{B,\diamond}}
  \DeltaPrem{Bijektivität}{f\colon A\bij B}
  \DeltaRow{\textbf{Neues Symbol}}{}
  \DeltaPrem{Isomorphismus}{%
    \IdempotentSemigroupIso{f}{A,\star}{B,\diamond}}
}

\FormulaAxiomDeltaK[Homomorphismuszugriff]{%
  \IdempotentSemigroupIso{f}{A,\star}{B,\diamond}
  \vdash\IdempotentSemigroupHom{f}{A,\star}{B,\diamond}%
}{IdempotentSemigroupIsoHomAxiom}{%
  \DeltaRow{Mengen}{A\dsep B}
  \DeltaRow{Funktionen}{f\dsep\star\dsep\diamond}
}

\FormulaAxiomDeltaK[Bijektivitätszugriff]{%
  \IdempotentSemigroupIso{f}{A,\star}{B,\diamond}
  \vdash f\colon A\bij B%
}{IdempotentSemigroupIsoBijectiveAxiom}{%
  \DeltaRow{Mengen}{A\dsep B}
  \DeltaRow{Funktionen}{f\dsep\star\dsep\diamond}
}

%============================================================
\chapter{Rekonstruktion aus der Potenzhalbgruppe}
%============================================================

In diesem Kapitel schreiben wir die Halbgruppenoperation als Juxtaposition
und kürzen \(\PowNE{S}\) durch \(\mathcal P_+(S)\) ab.  Für
\(a\in S\) steht \(\Psi(a)\) stets für \(\Psi(\{a\})\); diese Menge
muss nicht einelementig sein.  Gerade diese Unterscheidung ist wesentlich:
Bei Linksnullhalbgruppen können Isomorphismen der Potenzhalbgruppen
Einermengen auf größere Mengen abbilden.

Ziel ist die für die endliche Potenzhalbgruppenvermutung benötigte
Richtung
\[
  \mathcal P_+(S)\cong\mathcal P_+(T)
  \quad\Longrightarrow\quad S\cong T
\]
für endliche idempotente Halbgruppen.  Der Beweis folgt der
Rekonstruktionsidee von Yu--Zhao--Gan, trennt aber den endlichen
Kardinalitätsschritt ausdrücklich von den strukturellen Argumenten.
Bis zum Hauptsatz seien die betrachteten Grundhalbgruppen nichtleer; der
Hauptsatz behandelt den Leerfall vor Anwendung der Bandzerlegung gesondert.

\section{Rechtecksbänder und Bandzerlegung}

Eine idempotente Halbgruppe \(R\) heißt \emph{Rechtecksband}, wenn
\[
  xyx=x\qquad(x,y\in R).
\]
In einem Rechtecksband gilt äquivalent \(xyz=xz\).  Jedes
Rechtecksband besitzt eine Darstellung
\[
  R\cong I\times J,
  \qquad (i,j)(k,\ell)=(i,\ell),                         \tag{24.1}
\]
mit nichtleeren Mengen \(I,J\).  Der Fall \(|J|=1\) ist ein
Linksnullband, der Fall \(|I|=1\) ein Rechtsnullband; sind beide Faktoren
einelementig, so ist \(R\) trivial.  Sind \(|I|,|J|>1\), nennen wir
\(R\) ein \emph{echtes Rechtecksband}.

Wir verwenden den Struktursatz für Bänder: Jede idempotente Halbgruppe
besitzt eine maximale Halbverbandszerlegung
\[
  S=[Y;S_\alpha],\qquad
  S=\mathop{\dot\bigcup}_{\alpha\in Y}S_\alpha,           \tag{24.2}
\]
wobei \(Y\) eine kommutative idempotente Halbgruppe, jedes
\(S_\alpha\) ein Rechtecksband und
\[
  S_\alpha S_\beta\subseteq S_{\alpha\beta}              \tag{24.3}
\]
ist.  Die Ordnung auf \(Y\) ist durch
\(\beta\leq\alpha\iff\alpha\beta=\beta\) gegeben.
Konkret ist (24.2) der Quotient nach der kleinsten Kongruenz, deren
Quotient kommutativ ist; ihre Klassen sind genau die maximalen
Rechtecksbänder.  Idempotenz liefert die Idempotenz des Quotienten,
und (24.3) folgt aus der Kongruenzeigenschaft.  Dies ist die
McLean-Zerlegung eines Bandes.

Der folgende Komponentensatz ist der strukturelle Ausgangspunkt des
Rekonstruktionsbeweises.  Er ist wichtig genug, um seine beiden Aussagen
getrennt festzuhalten.

\begin{quote}
\textbf{Komponentensatz.}
Seien
\(S=[Y;S_\alpha]\) und \(T=[Z;T_\gamma]\) Bandzerlegungen und sei
\[
  \Psi\colon\mathcal P_+(S)\longrightarrow\mathcal P_+(T)
\]
ein Halbgruppenisomorphismus.  Dann existiert ein
Halbverbandsisomorphismus \(\theta\colon Y\to Z\) mit
\[
  \Psi\bigl[\mathcal P_+(S_\alpha)\bigr]
     =\mathcal P_+(T_{\theta(\alpha)})                    \tag{24.4}
\]
für jedes \(\alpha\in Y\).  Sind
\(a\in S_\alpha\), \(b\in S_\beta\) und
\(\alpha\nleq\beta\), so ist außerdem
\[
  \Psi(aba)                                                \tag{24.5}
\]
eine Einermenge.
\end{quote}

Die Komponentenerkennung (24.4) ist Satz~3.3 und die
Sandwicherkennung (24.5) ist Proposition~3.8(ii) in
\href{https://doi.org/10.1080/00927872.2015.1087006}%
  {Gan--Zhao--Shao, \emph{Globals of Idempotent Semigroups}, 2016}.
Die dort für \(\alpha>\beta\) formulierte Sandwichaussage liefert
(24.5) auch unter der schwächeren Voraussetzung
\(\alpha\nleq\beta\): Dann ist
\(c=aba\in S_{\alpha\beta}\), es gilt
\(\alpha>\alpha\beta\) und \(aca=c\), sodass die Aussage auf das Paar
\(a,c\) anwendbar ist.
Sie beruhen auf der maximalen Halbverbandszerlegung der Potenzhalbgruppe;
insbesondere setzen sie die zu rekonstruierenden Einermengen nicht voraus.
Damit entsteht kein Zirkelschluss mit dem folgenden Beweis.

\paragraph{Endliche echte Rechtecksbänder erkennen Einermengen.}
Seien \(R=I\times J\) und \(R'=I'\times J'\) endliche echte
Rechtecksbänder und
\(\Phi\colon\mathcal P_+(R)\to\mathcal P_+(R')\) ein
Isomorphismus.  Dann bildet \(\Phi\) jede Einermenge auf eine Einermenge
ab und induziert einen Isomorphismus \(R\to R'\).

\begin{proof}
Für nichtleere \(A,B\subseteq I\times J\) ergibt (24.1)
\[
  AB=\pi_I(A)\times\pi_J(B).                              \tag{24.6}
\]
Definiere rein multiplikativ
\[
\begin{aligned}
A\equiv_R B&\iff (\forall C\in\mathcal P_+(R))\ AC=BC,\\
A\equiv_L B&\iff (\forall C\in\mathcal P_+(R))\ CA=CB.
\end{aligned}
\]
Nach (24.6) gilt
\[
  A\equiv_RB\iff\pi_I(A)=\pi_I(B),\qquad
  A\equiv_LB\iff\pi_J(A)=\pi_J(B).                     \tag{24.7}
\]
  Hat \(K\subseteq I\) genau \(k\) Elemente, so enthält die zu
  \(K\) gehörende \(\equiv_R\)-Klasse genau
  \((2^{|J|}-1)^k\) Mengen: Eine Menge mit erster Projektion \(K\)
  wird eindeutig dadurch festgelegt, dass man für jedes \(i\in K\)
  eine nichtleere Teilmenge der \(i\)-ten Zeile wählt; umgekehrt
  liefert jede derartige Familie genau eine solche Menge.  Weil
\(|J|>1\), sind die kleinsten \(\equiv_R\)-Klassen genau diejenigen
mit \(k=1\).  Dual sind wegen \(|I|>1\) die kleinsten
\(\equiv_L\)-Klassen genau diejenigen mit einelementiger
\(J\)-Projektion.

Der Isomorphismus \(\Phi\) erhält beide Relationen in (24.7) und die
endlichen Größen ihrer Klassen.  Daher besitzt \(A\) genau dann auf
beiden Seiten einelementige Projektion, wenn \(\Phi(A)\) diese
Eigenschaft besitzt.  Eine nichtleere Teilmenge von \(I\times J\) mit
zwei einelementigen Projektionen ist eine Einermenge.  Also erhält
\(\Phi\) die Einermengen.  Die Einschränkung von \(\Phi\) auf ihnen
ist bijektiv und wegen
\(\Phi(\{x\}\{y\})=\Phi(\{xy\})\) multiplikativ.
\end{proof}

Ist eine Komponente ein Links- beziehungsweise Rechtsnullband, so ist auch
die durch (24.4) zugeordnete Komponente vom gleichen Typ: In ihrer
Potenzhalbgruppe gilt dann \(AB=A\) beziehungsweise \(AB=B\), und die
Gleichung für Einermengen überträgt sich auf die Grundkomponente.  Eine
triviale Komponente bleibt trivial, weil (24.4) eine Bijektion ihrer
nichtleeren Potenzmengen liefert.  Somit stimmen die vier Komponententypen
unter \(\theta\) überein.

\section{Transport innerhalb der Komponenten}

Von nun an seien \(S,T\) endlich, und \(\theta\) und \(\Psi\) seien wie
im Komponentensatz.  Für eine nichtleere Teilmenge \(A\subseteq S\) und
\(U\subseteq T\) setzen wir
\[
\begin{aligned}
 \widehat\Psi(A)
   &\coloneqq\bigcup_{\varnothing\neq C\subseteq A}\Psi(C),\\
 \widehat{\Psi^{-1}}(U)
   &\coloneqq\bigcup_{\varnothing\neq V\subseteq U}\Psi^{-1}(V).
\end{aligned}                                                   \tag{24.8}
\]
Beide Operatoren sind monoton.  Ferner gelten die Inklusionen
\[
 \Psi(A)\subseteq\widehat\Psi(A),
 \qquad
 A\subseteq\widehat{\Psi^{-1}}(\Psi(A)),
\]
sowie ihre dualen Entsprechungen.

\paragraph{Sandwichrechnung.}
Sind \(a\in S_\alpha\), \(b\in S_\beta\) und
\(\alpha\nleq\beta\), so gilt
\begin{align}
&\widehat{\Psi^{-1}}(\Psi(a))\,
 \widehat{\Psi^{-1}}(\Psi(b))\,
 \widehat{\Psi^{-1}}(\Psi(a))=\{aba\},                  \tag{24.9}\\
&\widehat\Psi\!\left(\widehat{\Psi^{-1}}(\Psi(a))\right)
 \widehat\Psi\!\left(\widehat{\Psi^{-1}}(\Psi(b))\right)
 \widehat\Psi\!\left(\widehat{\Psi^{-1}}(\Psi(a))\right)
 =\Psi(aba).                                               \tag{24.10}
\end{align}

\begin{proof}
Wähle Elemente der drei Faktoren auf der linken Seite von (24.9).  Nach
(24.8) stammen sie aus Urbildern nichtleerer Mengen
\(U_1,U_2\subseteq\Psi(a)\) und
\(V\subseteq\Psi(b)\).  Das Produkt \(U_1VU_2\) ist nichtleer und
liegt in
\[
  \Psi(a)\Psi(b)\Psi(a)=\Psi(aba).
\]
Die rechte Seite ist nach (24.5) eine Einermenge, also gilt
\(U_1VU_2=\Psi(aba)\).  Da \(\Psi^{-1}\) multiplikativ ist,
\[
  \Psi^{-1}(U_1)\Psi^{-1}(V)\Psi^{-1}(U_2)
  =\Psi^{-1}(U_1VU_2)=\{aba\}.
\]
Dies beweist (24.9), weil \(a,b,a\) selbst in den drei Faktoren liegen.
Wendet man \(\Psi\) auf die in (24.9) auftretenden nichtleeren
Teilmengen an und vereinigt wie in (24.8), so folgt (24.10).
\end{proof}

Aus (24.9) folgt zunächst noch folgende Mengenform.  Sind
\(a\in S_\alpha\), \(\alpha>\beta\),
\(B\in\mathcal P_+(S_\beta)\) und
\(A_1,A_2\) nichtleere Teilmengen von
\(\widehat{\Psi^{-1}}(\Psi(a))\), so gilt
\[
  A_1BA_2=aBa.                                             \tag{24.10c}
\]
Denn für jedes \(b\in B\) ist (24.9), auf die entsprechenden
Teilmengen eingeschränkt, genau
\(A_1\{b\}A_2=\{aba\}\); die Vereinigung über \(b\in B\) ergibt
(24.10c).

Insbesondere folgen für \(\alpha>\beta\) die beiden Formen
\[
 \Psi(a)t\Psi(a)\text{ ist einelementig}
 \quad(a\in S_\alpha,\ t\in T_{\theta(\beta)}),           \tag{24.10a}
\]
und
\[
 s\Psi(b)s\text{ ist einelementig}
 \quad(s\in T_{\theta(\alpha)},\ b\in S_\beta).          \tag{24.10b}
\]
Für (24.10a) wähle \(r\in\Psi(a)\) und setze
\(A_r=\Psi^{-1}(\{r\})\),
\(B_t=\Psi^{-1}(\{t\})\).  Dann ist
\(A_r\subseteq\widehat{\Psi^{-1}}(\Psi(a))\), und (24.10c) liefert
\[
 A_rB_tA_r=aB_ta.
\]
Nach Anwenden von \(\Psi\) ist daher
\(\Psi(a)t\Psi(a)=\{rtr\}\).  Für (24.10b) wenden wir (24.10a)
zunächst auf \(\Psi^{-1}\) an.  Dadurch ist
\(\Psi^{-1}(\{s\})\,b\,\Psi^{-1}(\{s\})\) einelementig, etwa gleich
\(\{cbc\}\) mit \(c\in\Psi^{-1}(\{s\})\).  Nach Anwenden von \(\Psi\)
ist \(s\Psi(b)s=\Psi(cbc)\), und diese Menge ist nach (24.5)
einelementig.

Für \(\alpha>\beta\) definieren wir auf \(S_\alpha\)
\[
 x\mathrel{\sigma_{\alpha,\beta}}y
 \iff
 (\forall b\in S_\beta)\ (xb=yb\ \land\ bx=by),          \tag{24.11}
\]
und für \(\gamma>\alpha\) auf \(S_\alpha\)
\[
 x\mathrel{\kappa_{\alpha,\gamma}}y
 \iff
 (\forall c\in S_\gamma)\ cxc=cyc.                      \tag{24.12}
\]
Beide Relationen sind Kongruenzen.  Für \(\sigma\) folgt dies unmittelbar
daraus, dass Multiplikation mit einem Element aus \(S_\alpha\) die
Testelemente in \(S_\beta\) belässt.  Sind nämlich
\(x\mathrel{\sigma_{\alpha,\beta}}y\),
\(d\in S_\alpha\) und \(b\in S_\beta\), so gilt
\[
\begin{aligned}
 (dx)b&=d(xb)=d(yb)=(dy)b,&
 b(dx)&=(bd)x=(bd)y=b(dy),\\
 (xd)b&=x(db)=y(db)=(yd)b,&
 b(xd)&=(bx)d=(by)d=b(yd),
\end{aligned}
\]
weil \(bd,db\in S_\beta\).  Für \(\kappa\) zeigt etwa
\[
 c(xd)c=cx(xc)dc=(cxc)dc=(cyc)dc=cy(yc)dc=c(yd)c
\]
für \(d\in S_\alpha\), \(c\in S_\gamma\) die Verträglichkeit mit
Rechtsmultiplikation; links gilt das duale Argument.  Dabei wurde in der
Rechteckskomponente \(S_\alpha\) die Identität \(uvw=uw\) benutzt.

Aus der Idempotenz erhält man außerdem die nützliche Charakterisierung
\[
 x\mathrel{\sigma_{\alpha,\beta}}y
 \iff
 (\forall b\in S_\beta)\ xbx=yby.                         \tag{24.13}
\]
Die Hinrichtung folgt aus
\(xbx=(xb)(bx)\).  Umgekehrt liefern
\(xb=(xbx)b\) und \(bx=b(xbx)\) die beiden Gleichungen in
(24.11).

Setze
\[
 \sigma_\alpha=\bigcap_{\beta<\alpha}\sigma_{\alpha,\beta},
 \qquad
 \kappa_\alpha=\bigcap_{\gamma>\alpha}\kappa_{\alpha,\gamma},
 \qquad
 \rho_\alpha=\sigma_\alpha\cap\kappa_\alpha,             \tag{24.14}
\]
wobei ein leerer Durchschnitt die Allrelation auf \(S_\alpha\) bedeutet.
Dann ist \(\rho_\alpha\) eine Kongruenz.  Auf den Komponenten von \(T\)
werden \(\sigma',\kappa',\rho'\) gleichartig definiert.

\begin{quote}
\textbf{Transportlemma.}
Für \(a,b\in S_\alpha\) gilt
\[
 a\mathrel{\rho_\alpha}b
 \iff
 (\forall s,t\in\Psi(a)\cup\Psi(b))\;
 s\mathrel{\rho'_{\theta(\alpha)}}t.                    \tag{24.15}
\]
Ist \(C=a\rho_\alpha\) und \(t\in\Psi(a)\), so bildet \(\Psi\)
\(\mathcal P_+(C)\) isomorph auf
\(\mathcal P_+(t\rho'_{\theta(\alpha)})\) ab.  Die Vorschrift
\[
 \tau_\alpha(a\rho_\alpha)
   =t\rho'_{\theta(\alpha)}\qquad(t\in\Psi(a))            \tag{24.16}
\]
definiert daher einen Isomorphismus
\[
 \tau_\alpha\colon S_\alpha/\rho_\alpha
    \longrightarrow T_{\theta(\alpha)}/\rho'_{\theta(\alpha)}.
                                                                  \tag{24.17}
\]
Schließlich gelten für \(a\in S_\alpha\), \(b\in S_\beta\):
\begin{align}
 \alpha,\beta>\alpha\beta:\quad&
 \tau_\alpha(a\rho_\alpha)\,
 \tau_\beta(b\rho_\beta)=\Psi(ab),                       \tag{24.18}\\[2mm]
 \alpha>\beta:\quad&
 \tau_\alpha(a\rho_\alpha)\,
 \tau_\beta(b\rho_\beta)\,
 \tau_\alpha(a\rho_\alpha)=\Psi(aba).                    \tag{24.19}
\end{align}
Die Produkte auf den linken Seiten sind jeweils Einermengen.
\end{quote}

\begin{proof}
Wir geben die Transportrechnung ausführlich an.

Sei zunächst \(x\mathrel{\sigma_{\alpha,\beta}}y\).  Nach (24.13)
und der Sandwichrechnung sind alle Elemente einer Menge \(\Psi(x)\)
untereinander \(\sigma'_{\theta(\alpha),\theta(\beta)}\)-äquivalent.
Für \(B'\in\mathcal P_+(T_{\theta(\beta)})\) liegt wegen (24.4)
\(B=\Psi^{-1}(B')\) in \(\mathcal P_+(S_\beta)\).  Aus (24.11) folgt
\(xB=yB\) und \(Bx=By\), also
\[
 \Psi(x)B'=\Psi(y)B',\qquad
 B'\Psi(x)=B'\Psi(y).
\]
Zusammen mit der bereits erhaltenen Äquivalenz innerhalb jedes Bildes
zeigt dies
\[
 s\mathrel{\sigma'_{\theta(\alpha),\theta(\beta)}}t
 \quad(s,t\in\Psi(x)\cup\Psi(y)).                         \tag{24.20}
\]

Sei nun \(x\mathrel{\kappa_{\alpha,\gamma}}y\) und
\(c'\in T_{\theta(\gamma)}\).  Das Urbild
\(A=\Psi^{-1}(\{c'\})\) liegt nach (24.4) in
  \(\mathcal P_+(S_\gamma)\).  Wähle \(c\in A\).  Wendet man
  (24.10a) auf den inversen Isomorphismus \(\Psi^{-1}\) an, so sind
  \(AxA\) und \(AyA\) Einermengen.  Da sie \(cxc\) beziehungsweise
  \(cyc\) enthalten, liefert (24.12)
\[
 AxA=\{cxc\}=\{cyc\}=AyA.
\]
Nach Anwenden von \(\Psi\) ist
\(c'\Psi(x)c'=c'\Psi(y)c'\); diese Menge ist nach (24.10b)
einelementig.  Folglich gilt
\[
 s\mathrel{\kappa'_{\theta(\alpha),\theta(\gamma)}}t
 \quad(s,t\in\Psi(x)\cup\Psi(y)).                       \tag{24.21}
\]

Wir benötigen auch die Mengenfassung dieser beiden Aussagen.  Ist
\(D\subseteq a\sigma_{\alpha,\beta}\) nichtleer und
\(u\in T_{\theta(\beta)}\), so ist
\(B=\Psi^{-1}(\{u\})\subseteq S_\beta\) und \(B^2=B\), denn
\(u^2=u\) und \(\Psi^{-1}\) ist multiplikativ.  Aus
\(DB=aB\) und \(BD=Ba\) folgt
\[
 DBD=aBa.
\]
Nach Anwenden von \(\Psi\) ist deshalb
\(\Psi(D)u\Psi(D)=\Psi(a)u\Psi(a)\).  Die rechte Seite ist nach der
Sandwichrechnung einelementig; (24.13) zeigt somit
\[
 s\mathrel{\sigma'_{\theta(\alpha),\theta(\beta)}}t
 \quad(s,t\in\Psi(a)\cup\Psi(D)).                         \tag{24.20a}
\]
Ist entsprechend \(D\subseteq a\kappa_{\alpha,\gamma}\), setze
\(A=\Psi^{-1}(\{v\})\subseteq S_\gamma\) für
\(v\in T_{\theta(\gamma)}\).  Die zu (24.10a) symmetrische Aussage für
den Isomorphismus \(\Psi^{-1}\) zeigt, dass \(AdA\) für jedes
\(d\in D\cup\{a\}\) einelementig ist.  Für ein festes
\(c\in A\) gilt somit \(AdA=\{cdc\}\).  Aus (24.12) folgt
\(cdc=cac\), also \(ADA=AaA\).  Nach Anwenden von \(\Psi\) gilt
\(v\Psi(D)v=v\Psi(a)v\).  Diese gemeinsame Menge ist nach (24.10b)
einelementig; die Definition (24.12) ergibt daher
\[
 s\mathrel{\kappa'_{\theta(\alpha),\theta(\gamma)}}t
 \quad(s,t\in\Psi(a)\cup\Psi(D)).                        \tag{24.21a}
\]

Die Schnitte in (24.14) liefern aus (24.20)--(24.21) die Hinrichtung
von (24.15).  Für die Rückrichtung wähle
\(s\in\Psi(a)\) und setze \(A_s=\Psi^{-1}(\{s\})\).  Die Voraussetzung
legt sowohl \(\Psi(a)\) als auch \(\Psi(b)\) in die
\(\rho'_{\theta(\alpha)}\)-Klasse von \(s\).  Wendet man
(24.20a)--(24.21a) für den inversen Isomorphismus zuerst auf
\(\Psi(a)\) und dann auf \(\Psi(b)\) an, so liegen alle Elemente von
\[
 A_s\cup\Psi^{-1}(\Psi(a))=A_s\cup\{a\}
 \quad\text{beziehungsweise}\quad
 A_s\cup\Psi^{-1}(\Psi(b))=A_s\cup\{b\}
\]
jeweils in einer \(\rho_\alpha\)-Klasse.  Die nichtleere Brückenmenge
\(A_s\) zeigt daher \(a\mathrel{\rho_\alpha}b\).

Damit liegt nach (24.20a)--(24.21a) für jede nichtleere Teilmenge
\(D\subseteq a\rho_\alpha\) das Bild \(\Psi(D)\) vollständig in der Klasse
\(t\rho'_{\theta(\alpha)}\).  Sei umgekehrt
\(E\subseteq t\rho'_{\theta(\alpha)}\) nichtleer und
\(D=\Psi^{-1}(E)\).  Die mengenweisen Transportaussagen für
\(\Psi^{-1}\), einmal mit \(E\) und einmal mit \(\Psi(a)\), zeigen,
dass sowohl \(D\) als auch \(\{a\}=\Psi^{-1}(\Psi(a))\) in derselben
\(\rho_\alpha\)-Klasse wie die nichtleere Brückenmenge
\(\Psi^{-1}(\{t\})\) liegen.  Also ist
\(D\subseteq a\rho_\alpha\).  Die behauptete Einschränkung von \(\Psi\)
ist somit bijektiv und multiplikativ.

Die Wohldefiniertheit und Injektivität von (24.16) folgen aus (24.15).
Für die Surjektivität sei \(E=t\rho'_{\theta(\alpha)}\) eine beliebige
Zielklasse.  Wähle \(a\in\Psi^{-1}(\{t\})\).  Die soeben bewiesene
Klassenrestriktion, nun für den inversen Isomorphismus, bildet
\(\mathcal P_+(E)\) bijektiv auf
\(\mathcal P_+(a\rho_\alpha)\) ab.  Äquivalent bildet \(\Psi\) die
zweite Potenzhalbgruppe auf die erste ab; insbesondere liegt
\(\Psi(a)\) in \(E\), und (24.16) bildet \(a\rho_\alpha\) auf \(E\).
Damit ist \(\tau_\alpha\) bijektiv.  Sind
\(a,b\in S_\alpha\), \(s\in\Psi(a)\) und
\(t\in\Psi(b)\), so liegt \(st\) in \(\Psi(ab)\); daher
\[
 \tau_\alpha(a\rho_\alpha)\tau_\alpha(b\rho_\alpha)
  =(st)\rho'_{\theta(\alpha)}
  =\tau_\alpha((ab)\rho_\alpha).
\]
Somit ist (24.17) ein Isomorphismus.

Für (24.18) setzen wir \(\delta=\alpha\beta\).  Aus
\(x\mathrel{\sigma_{\alpha,\delta}}a\) und
\(y\mathrel{\sigma_{\beta,\delta}}b\) folgt durch die Kette
\[
\begin{aligned}
xy&=(xyx)y=(xyx)b=(xy)ab\\
  &=xy(ab)=xb(ab)=a(bab)=ab
\end{aligned}                                               \tag{24.22}
\]
für alle \(x\in a\rho_\alpha\), \(y\in b\rho_\beta\);
jeweils wurde (24.11) benutzt, und alle eingeklammerten Testelemente
liegen in \(S_\delta\).  Die beiden \(\rho'\)-Klassen auf der linken
Seite von (24.18) liegen nach (24.14) in den entsprechenden
\(\sigma'\)-Klassen.  Daher gilt die Rechnung (24.22) ebenso für
beliebige Vertreter dieser Zielklassen, und das linke Produkt in (24.18)
kollabiert auf ein Element.  Es enthält \(st\) für
\(s\in\Psi(a)\), \(t\in\Psi(b)\), und dieses Element liegt in
\(\Psi(ab)\).  Da auch \(\Psi(ab)\) durch dieselbe Rechnung nur dieses
Element enthält, folgt (24.18).

Für (24.19) seien \(x\in a\rho_\alpha\) und
\(y\in b\rho_\beta\).  Dann gilt
\(x\mathrel{\sigma_{\alpha,\beta}}a\) und
\(y\mathrel{\kappa_{\beta,\alpha}}b\).  Mit (24.13) und (24.12)
folgt
\[
  xyx=aya=aba.                                             \tag{24.23}
\]
Nach (24.14) liegen die entsprechenden Zielklassen in den benötigten
\(\sigma'\)- und \(\kappa'\)-Klassen; daher gilt (24.23) auch für alle
ihre Vertreter.  Das linke Produkt von (24.19) ist somit einelementig; es
enthält \(sts\in\Psi(aba)\) für
\(s\in\Psi(a)\), \(t\in\Psi(b)\).  Gleichung (24.5) macht
\(\Psi(aba)\) selbst einelementig, und (24.19) folgt.
\end{proof}

\section{Endliches Anheben auf die Grundelemente}

Der Quotientenisomorphismus \(\tau_\alpha\) muss nun innerhalb jeder
\(\rho_\alpha\)-Klasse zu einer Bijektion der Grundelemente angehoben
werden.  Bei echten Rechtecksbändern und bei trivialen Komponenten ist
diese Hebung bereits kanonisch.  Für echte Rechtecksbänder folgt die
Einermengeneigenschaft aus dem oben bewiesenen Erkennungssatz; bei einer
trivialen Komponente folgt sie unmittelbar aus (24.4), weil ihre
nichtleere Potenzmenge genau ein Element besitzt.  In beiden Fällen ist
\(\Psi(a)\) für jedes \(a\in S_\alpha\) eine Einermenge, und wir setzen
\[
  \eta_\alpha(a)=\text{das einzige Element von }\Psi(a).   \tag{24.24}
\]
Die so definierte Abbildung ist ein Isomorphismus
\(S_\alpha\to T_{\theta(\alpha)}\).  Wegen (24.16) liegt
\(\eta_\alpha(a)\) außerdem in der Klasse
\(\tau_\alpha(a\rho_\alpha)\).

Es bleiben Links- und Rechtsnullkomponenten.  Für eine solche Komponente
setzen wir
\[
U_\alpha=\begin{cases}
\displaystyle
 \bigcup_{\gamma>\alpha}\ \bigcup_{c\in S_\gamma}cS_\alpha,
   &\text{falls }S_\alpha\text{ linksnull ist},\\[3mm]
\displaystyle
 \bigcup_{\gamma>\alpha}\ \bigcup_{c\in S_\gamma}S_\alpha c,
   &\text{falls }S_\alpha\text{ rechtsnull ist}.
\end{cases}                                                \tag{24.25}
\]
Für eine zugeordnete Komponente von \(T\) bezeichnet
\(U_{\theta(\alpha)}\) im Folgenden die durch dieselbe Vorschrift in
\(T\) gebildete Menge.
Ist \(\alpha\) maximal, so ist \(U_\alpha=\varnothing\).  In beiden
Fällen gilt auch
\[
 U_\alpha=
 \bigcup_{\gamma>\alpha}\ \bigcup_{c\in S_\gamma}cS_\alpha c.  \tag{24.26}
\]
  Tatsächlich liegen \(cd\) und \(dc\) beide in \(S_\alpha\).  Ist diese
  Komponente linksnull, so ist
  \(cd=(cd)(dc)=cdc\); ist sie rechtsnull, so ist
  \(dc=(cd)(dc)=cdc\).  Dies gilt für
  \(c\in S_\gamma\), \(d\in S_\alpha\) und \(\gamma>\alpha\) und beweist
  (24.26).

\paragraph{Kanonischer Teil der Hebung.}
Für \(u\in U_\alpha\) ist \(\Psi(u)\) eine Einermenge, und ihr einziges
Element liegt in \(U_{\theta(\alpha)}\).  Die dadurch definierte Abbildung
\[
 U_\alpha\longrightarrow U_{\theta(\alpha)}              \tag{24.27}
\]
ist bijektiv.

\begin{proof}
Nach (24.26) hat \(u\) die Form \(cdc\) mit
\(c\in S_\gamma\), \(d\in S_\alpha\) und \(\gamma>\alpha\).
Die Sandwicherkennung (24.5) zeigt, dass
\[
 \Psi(u)=\Psi(c)\Psi(d)\Psi(c)
\]
einelementig ist.  Ihr Element hat die Form \(sts\) mit
\(s\in T_{\theta(\gamma)}\) und
\(t\in T_{\theta(\alpha)}\); nach (24.26) liegt es in
\(U_{\theta(\alpha)}\).  Die gleiche Aussage für \(\Psi^{-1}\) liefert
Surjektivität und zeigt zugleich, dass die beiden Zuordnungen invers sind.
\end{proof}

Sei jetzt \(C=a\rho_\alpha\) und
\[
 D=\tau_\alpha(C)
   =t\rho'_{\theta(\alpha)}\qquad(t\in\Psi(a)).            \tag{24.28}
\]
Nach dem Transportlemma schränkt sich \(\Psi\) zu einer Bijektion
\(\mathcal P_+(C)\to\mathcal P_+(D)\) ein.  Also sind die beiden
nichtleeren Potenzmengen gleichmächtig.  Weil \(C,D\) endlich sind,
liefert
\FormulaRefAuto{FiniteNonemptyPowerSetReflectsCardinality}[thm]
\[
  |C|=|D|.                                                  \tag{24.29}
\]
Aus (24.15) und (24.27) folgt außerdem, dass der kanonische Teil eine
Bijektion
\[
 C\cap U_\alpha\longrightarrow
 D\cap U_{\theta(\alpha)}                                 \tag{24.30}
\]
ist: Die Inklusion folgt aus (24.15), und die umgekehrte Inklusion folgt
durch dieselbe Argumentation für die inverse kanonische Abbildung.  Aus
(24.29) und (24.30) folgt, unter Verwendung der disjunkten Zerlegungen
\(C=(C\cap U_\alpha)\cup(C\setminus U_\alpha)\) und der
analogen Zerlegung von \(D\),
\[
\begin{aligned}
 |C\setminus U_\alpha|
 &=|C|-|C\cap U_\alpha|\\
 &=|D|-|D\cap U_{\theta(\alpha)}|\\
 &=|D\setminus U_{\theta(\alpha)}|.
\end{aligned}
\]
Ergänze (24.30) auf diesen Komplementen durch eine beliebige Bijektion.
So entsteht eine Bijektion
\[
  \overline\tau_C\colon C\longrightarrow D               \tag{24.31}
\]
mit
\[
 \overline\tau_C(u)=\text{das einzige Element von }\Psi(u)
 \qquad(u\in C\cap U_\alpha).                             \tag{24.32}
\]
Da \(\rho_\alpha\) und \(\rho'_{\theta(\alpha)}\) Kongruenzen sind und
die Grundelemente idempotent sind, sind \(C\) und \(D\) Unterhalbgruppen.
Jede Abbildung (24.31) ist ein Isomorphismus, denn Quelle und Ziel sind
beide Linksnull- beziehungsweise beide Rechtsnullhalbgruppen.

Die Menge der \(\rho_\alpha\)-Klassen ist das Bild der endlichen Menge
\(S_\alpha\) unter der Quotientenprojektion und daher endlich.  Durch
endlich viele aufeinanderfolgende Wahlen wählen wir für jede Klasse genau
eine solche Abbildung und vereinigen sie; wegen der paarweise disjunkten
Klassen ist diese Vereinigung eine Funktion.  Man erhält einen
Komponentenisomorphismus
\[
 \eta_\alpha\colon S_\alpha\longrightarrow
 T_{\theta(\alpha)}                                       \tag{24.33}
\]
mit den beiden Verträglichkeiten
\begin{align}
 \eta_\alpha(a)\rho'_{\theta(\alpha)}
    &=\tau_\alpha(a\rho_\alpha),                         \tag{24.34}\\
 a\in U_\alpha\quad&\Longrightarrow\quad
 \eta_\alpha(a)=\text{das einzige Element von }\Psi(a). \tag{24.35}
\end{align}
Dabei gilt (24.34) auch für die zuvor kanonisch definierten Abbildungen
aus (24.24).  Damit sind die Abbildungen \(\eta_\alpha\) für alle vier
Komponententypen definiert.

\section{Globale Bestimmtheit endlicher Bänder}

\FormulaThmDeltaK[Endliche idempotente Halbgruppen sind durch ihre
Potenzhalbgruppen global bestimmt]{%
  \begin{aligned}[t]
   &\IdempotentSemigroupAlg{A}{\star}\dsep
    \IdempotentSemigroupAlg{B}{\diamond}\dsep
    \Finite{A}\dsep\Finite{B}\dsep{}\\[-2pt]
   &\SemigroupIso{\Psi}
      {\PowNE{A},\star_{\mathcal P}}
      {\PowNE{B},\diamond_{\mathcal P}}
    \vdash
    \exists f\,\IdempotentSemigroupIso{f}{A,\star}{B,\diamond}
  \end{aligned}%
}{FiniteIdempotentSemigroupsGloballyDetermined}{%
  \DeltaRow{Mengen}{A\dsep B\dsep X\dsep Y}
  \DeltaRow{Funktionen}{f\dsep\Psi\dsep\star\dsep\diamond\dsep
    \star_{\mathcal P}\dsep\diamond_{\mathcal P}}
}

\begin{proof}
Wir verwenden für \((A,\star)\) und \((B,\diamond)\) die Bezeichnungen
\(S\) und \(T\).  Ist \(A=\varnothing\), so ist
\(\PowNE{A}=\varnothing\).  Die Bijektivität von \(\Psi\) erzwingt
\(\PowNE{B}=\varnothing\) und damit \(B=\varnothing\).  Die leere
Abbildung ist dann der gewünschte Isomorphismus.  Ist umgekehrt
\(A\neq\varnothing\), so folgt aus der Bijektivität ebenso
\(B\neq\varnothing\).

Im verbleibenden nichtleeren Fall verwenden wir die Zerlegungen
\(S=[Y;S_\alpha]\) und \(T=[Z;T_\gamma]\).  Der Komponentensatz liefert
\(\theta\colon Y\to Z\), und die vorangegangene Konstruktion liefert für
jedes \(\alpha\in Y\) einen Isomorphismus
\(\eta_\alpha\colon S_\alpha\to T_{\theta(\alpha)}\).
Setze
\[
  \eta=\bigcup_{\alpha\in Y}\eta_\alpha.                 \tag{24.36}
\]
Wegen der disjunkten Zerlegungen, der Bijektivität von \(\theta\) und
der komponentenweisen Bijektivität ist \(\eta\colon S\to T\) bijektiv.
Es bleibt die Multiplikativität.

Seien \(a\in S_\alpha\) und \(b\in S_\beta\).

\smallskip
\noindent\emph{Erster Fall: \(\alpha=\beta\).}
Dann ist \(ab\in S_\alpha\), und (24.33) ergibt
\[
  \eta(ab)=\eta_\alpha(ab)
  =\eta_\alpha(a)\eta_\alpha(b)=\eta(a)\eta(b).          \tag{24.37}
\]

\smallskip
\noindent\emph{Zweiter Fall:
\(\delta=\alpha\beta\notin\{\alpha,\beta\}\).}
Dann gilt \(\alpha,\beta>\delta\) und \(ab\in S_\delta\).  Ist
\(S_\delta\) trivial oder ein echtes Rechtecksband, so folgt aus (24.24)
\[
  \eta(ab)=\text{das einzige Element von }\Psi(ab).
\]
Ist \(S_\delta\) ein Linksnullband, so gilt
\(ab=a(ab)\in S_\alpha S_\delta\subseteq U_\delta\); ist es ein
Rechtsnullband, so gilt
\(ab=(ab)b\in S_\delta S_\beta\subseteq U_\delta\).  In beiden Fällen
liefert (24.35) dieselbe Gleichheit.  Andererseits liegen
\(\eta(a)\) und \(\eta(b)\) nach (24.34) in den beiden Klassen auf der
linken Seite von (24.18).  Deshalb ist
\[
 \eta(a)\eta(b)
   =\text{das einzige Element von }\Psi(ab)
   =\eta(ab).                                               \tag{24.38}
\]

\smallskip
\noindent\emph{Dritter Fall:
\(\alpha\beta\in\{\alpha,\beta\}\) und \(\alpha\neq\beta\).}
Wir behandeln zunächst
\(\alpha=\alpha\beta<\beta\).  Zunächst sei \(S_\alpha\) trivial
oder ein echtes Rechtecksband.  Dann sind
\(\Psi(a)=\{\eta(a)\}\) und \(\Psi(ab)=\{\eta(ab)\}\).
Nach der Klassenbijektion des Transportlemmas existiert eine nichtleere
Menge \(C\subseteq b\rho_\beta\) mit
\(\Psi(C)=\{\eta(b)\}\).  Für jedes
\(s\in\Psi(b)\) gilt wegen
\(C\subseteq b\sigma_{\beta,\alpha}\)
\[
 \Psi(a)\Psi(C)=\Psi(a)\{s\}=\Psi(a)\Psi(b).              \tag{24.39}
\]
Nach (24.20a) sind \(\eta(b)\), das einzige Element von \(\Psi(C)\),
und jedes \(s\in\Psi(b)\) modulo
\(\sigma'_{\theta(\beta),\theta(\alpha)}\) äquivalent.  Die Definition
(24.11), nun in den Zielkomponenten, liefert die erste Gleichheit in
(24.39); die zweite folgt, weil sie für jedes \(s\in\Psi(b)\) gilt.
Somit folgt
\[
 \eta(a)\eta(b)=\Psi(a)\Psi(C)
   =\Psi(ab)=\eta(ab).                                     \tag{24.40}
\]

Ist \(S_\alpha\) linksnull, so liegt \(ab\) in \(S_\alpha\) und
\[
  ab=a(ab)=a.
\]
Auch \(T_{\theta(\alpha)}\) ist linksnull.  Mit
\(x=\eta(a)\eta(b)\in T_{\theta(\alpha)}\) gilt einerseits
\(\eta(a)x=\eta(a)\) durch die Linksnullidentität und andererseits
\(\eta(a)x=\eta(a)^2\eta(b)=x\).  Daher
\[
  \eta(a)\eta(b)=\eta(a)=\eta(ab).                        \tag{24.41}
\]

Ist \(S_\alpha\) rechtsnull, so liegen \(ba,ab\) in
\(S_\alpha\), und die Rechtsnullidentität ergibt
\[
  ab=(ba)(ab)=bab\in U_\alpha.                            \tag{24.42}
\]
Nach (24.35) ist \(\eta(ab)\) das einzige Element von \(\Psi(bab)\).
Im Ziel liegen \(x=\eta(a)\eta(b)\) und
\(y=\eta(b)\eta(a)\) beide im rechten Nullband
\(T_{\theta(\alpha)}\).  Daher ist \(yx=x\), also
\[
 \eta(a)\eta(b)=\eta(b)\eta(a)\eta(b).                   \tag{24.43}
\]
Wegen (24.34) liegt die rechte Seite in
\[
 \tau_\beta(b\rho_\beta)\,
 \tau_\alpha(a\rho_\alpha)\,
 \tau_\beta(b\rho_\beta).
\]
Nach (24.19), angewandt auf \(\beta>\alpha\), ist dieses Produkt genau
\(\Psi(bab)\).  Zusammen mit (24.42)--(24.43) folgt wieder
\(\eta(a)\eta(b)=\eta(ab)\).  Gilt stattdessen
\(\beta=\alpha\beta<\alpha\), so betrachten wir die
Oppositionshalbgruppen \(S^{\mathrm{op}}\) und \(T^{\mathrm{op}}\).
Dieselbe Abbildung \(\Psi\) ist auch ein Isomorphismus ihrer
Potenzhalbgruppen, denn die Multiplikation beider Potenzhalbgruppen wird
zugleich umgekehrt.  Dabei bleiben \(\sigma,\kappa,\rho\) und \(\tau\)
unverändert; Links- und Rechtsnullkomponenten tauschen ihre Rollen, und
die Vorschrift (24.25) liefert dieselben Mengen \(U_\alpha\).  Somit sind
auch dieselben Abbildungen \(\eta_\alpha\) zulässige Hebungen.  Auf das
Paar \((b,a)\) im Oppositum ist der bereits bewiesene Unterfall
anwendbar.  Dort sind
\(b\mathbin{\circ}a=ab\) und
\(\eta(b)\mathbin{\circ}\eta(a)=\eta(a)\eta(b)\); daher folgt die
gewünschte Gleichung auch in diesem Fall.

Damit ist \(\eta\) für alle \(a,b\in S\) multiplikativ.  Die bijektive
Abbildung \(\eta\) ist also ein Halbgruppenisomorphismus.  Da Quelle und
Ziel idempotent sind, ist sie nach der Definition zugleich ein
Isomorphismus idempotenter Halbgruppen.
\end{proof}

\section{Reichweite der Aussage}

\begin{remark}
Der Satz setzt auf beiden Seiten Idempotenz voraus.  Der Beweis erkennt
nicht aus \(\mathcal P_+(S)\cong\mathcal P_+(T)\), dass eine zunächst
beliebige Zielhalbgruppe \(T\) ebenfalls idempotent ist.  Eine solche
einseitige Erkennung wäre ein zusätzlicher und für die allgemeine
Potenzhalbgruppenvermutung wesentlicher Schritt.

Die in
\href{https://doi.org/10.1080/00927872.2017.1319474}%
 {Yu--Zhao--Gan, \emph{Global determinism of idempotent semigroups}, 2018}
ohne Endlichkeitsbedingung formulierte Aussage wird hier bewusst nicht
übernommen.  Für ein Linksnullband \(L_X\) ist
\(\mathcal P_+(L_X)\) selbst ein Linksnullband mit
\(2^{|X|}-1\) Elementen.  Bei unendlichen Kardinalzahlen folgt aus
\(2^\kappa=2^\lambda\) in ZFC nicht \(\kappa=\lambda\); entsprechende
Kontinuumsfunktionen sind nach
\href{https://doi.org/10.1016/0003-4843(70)90012-4}%
 {Eastons Satz}
mit ZFC verträglich.  Im endlichen Fall schließt dagegen
\FormulaRefAuto{FiniteNonemptyPowerSetReflectsCardinality}[thm]
genau diese Lücke.
\end{remark}

\end{document}
